% Tableau 3 — Coût-Bénéfice des mesures d'adaptation
% Sources : Alfieri et al. (2016); CE (2014) Guide CBA
\begin{table}[htbp]
\centering\small
\caption{Analyse co\^ut-b\'en\'efice des mesures d'adaptation (Alfieri et al. 2016; taux actualisation CE 2014)}
\label{tab:adaptation}
\begin{tabular}{p{3cm}rrlrrrl}
\hline
\textbf{Mesure} & \textbf{Co\^ut} & \textbf{Dur\'ee} & \textbf{Réd. RCP4.5} & \textbf{VAN\,3\%} & \textbf{VAN\,5\%} & \textbf{B/C} & \textbf{Note} \\
 & (M\euro) & (ans) & (\%) & (M\euro) & (M\euro) & & \\
\hline
Renforcement digues Meuse & 2,500 & 50 & 35.0% & 10,018 & 7,108 & 8.8× & \footnotesize{Efficacité réduite à RCP 8.5 : les } \\
Systèmes d'alerte précoce & 150 & 20 & 20.0% & 3,624 & 3,035 & 33.5× & \footnotesize{Rapport bénéfice/coût le plus élevé} \\
Reforestation bassins versan & 800 & 100 & 15.0% & 5,697 & 3,579 & 23.5× & \footnotesize{Ordre de grandeur pan-européen, inc} \\
Zones d'expansion inondables & 1,200 & 75 & 25.0% & 8,846 & 5,803 & 19.6× & \footnotesize{Implique des relocalisations — coût} \\
\hline
\end{tabular}
\end{table}
